\documentclass[titlepage]{article}

\usepackage[czech]{babel}
\usepackage{fontspec}
\usepackage{amsmath}
\usepackage{graphicx}
\usepackage{geometry}
\usepackage{hyperref}

\geometry{a4paper, margin=1in}

\def\MyTitle{Reziduální aritmetika}
\def\MySubtitle{Dokumentace k zápočtovému programu na P2X}
\def\MyAuthor{Ondřej Sedláček}
\def\MyDate{\today}

\newtheorem{theorem}{Věta}

\begin{document}

\begin{titlepage}
  \vbox to \textheight{
    \vfill
    
    \centering
    
    {\Huge\bfseries \MyTitle\par} 

    {\large\bfseries \MySubtitle\par}
    
    \vspace{2cm}
    
    {\LARGE \MyAuthor\par}
    
    \vspace{1cm}
    
    {\large \MyDate\par}

    \vfill
  }
\end{titlepage}
\clearpage
\pagenumbering{roman}

\tableofcontents
\listoffigures

\clearpage
\pagenumbering{arabic}

\section{Úvod}

K zápočtovému programu jsem si vybral téma Reziduální aritmetika, jehož součástí je implementace provádění aritmetických operací s dlouhými čísly za pomocí poznatků o zbytcích po dělení. Každé číslo v určitém rozmezí umíme jednoznačně vyjádřit jako seznam zbytků po dělení několika nesoudělnými čísly a s tímto seznamem pak velmi efektivně sčítat, odčítat a násobit.

Vytvořil jsem tedy knihovnu v jazyce C, která implementuje veškeré potřebné operace, jako je převod z desítkového zápisu na zápis se zbytky a naopak a pak operace sčítání, odčítání, násobení a porovnávání.

\section{Teorie}

Celý tento program stojí na poznatku z teorie čísel, které se nazývá Čínská zbytková věta. Níže je její znění:

\begin{theorem}[Čínská zbytková věta]
Nechť $n_1, n_2, \dots, n_k$ jsou po dvou nesoudělná přirozená čísla, tj. $\gcd(n_i, n_j) = 1$ pro $i \neq j$. Označme $N = n_1 n_2 \dots n_k$. Potom pro libovolná celá čísla $a_1, a_2, \dots, a_k$ má soustava kongruencí
$$
\begin{aligned}
    x &\equiv a_1 \pmod{n_1} \\
    x &\equiv a_2 \pmod{n_2} \\
      &\vdots \\
    x &\equiv a_k \pmod{n_k}
\end{aligned}
$$
řešení, které je určeno jednoznačně modulo $N$.
\end{theorem}



\end{document}
